\documentclass[a4paper,11pt]{article}
\usepackage{exam-style-341}
\examinehead{341 农业综合知识三 · 网络技术分章练习}{第02章\quad 协议与体系结构}
\begin{document}

\unitcover{341 农业综合知识三 · 网络技术分章练习}{第02章\quad 协议与体系结构}{练习（共 100 分）}{约 60 分钟}{分层思想与网络协议三要素；OSI 七层参考模型；TCP/IP 四层体系结构；五层原理体系结构；各层功能、PDU 与典型设备；封装与解封装；OSI 与 TCP/IP 的比较；客户／服务器方式与 P2P 方式；互联网标准化组织与 RFC。}

\begin{knowmap}
\textbf{分层思想：}把复杂的网络通信问题分解为若干层，每层完成特定功能、向上层提供服务并使用下层服务；好处是各层独立、灵活、易于实现和维护\\[2pt]
\textbf{协议三要素：}语法（数据与控制信息的格式）· 语义（发出何种控制信息、完成何种动作、返回何种应答）· 同步（事件实现顺序）\\[2pt]
\textbf{OSI 七层（自下而上）：}物理层 · 数据链路层 · 网络层 · 运输层 · 会话层 · 表示层 · 应用层\\[2pt]
\textbf{TCP/IP 四层：}网络接口层 · 网际层（IP）· 运输层（TCP、UDP）· 应用层\\[2pt]
\textbf{五层原理体系结构：}物理层 · 数据链路层 · 网络层 · 运输层 · 应用层（综合 OSI 与 TCP/IP 的优点）\\[2pt]
\textbf{封装与解封装：}发送时数据逐层加首部（应用层报文 $\rightarrow$ 运输层报文段 $\rightarrow$ 网络层分组 $\rightarrow$ 数据链路层帧 $\rightarrow$ 物理层比特流），接收时逐层剥去首部\\[2pt]
\textbf{通信方式：}客户／服务器方式（客户主动请求、服务器被动响应）与 P2P 方式（对等通信）
\end{knowmap}

\begin{keyterms}
\textbf{网络协议}：为进行网络中的数据交换而建立的规则、标准或约定，简称协议。\\
\textbf{语法}：规定数据与控制信息的格式，如帧的字段排列。\\
\textbf{语义}：规定需要发出何种控制信息、完成何种动作以及做出何种应答。\\
\textbf{同步}：规定各种事件实现顺序的详细说明，也称时序。\\
\textbf{网络体系结构}：计算机网络各层及其协议的集合，即网络的分层结构与各层协议的统称。\\
\textbf{OSI/RM}：国际标准化组织（ISO）提出的开放系统互连参考模型，共 7 层。\\
\textbf{TCP/IP 体系结构}：互联网实际使用的四层体系结构：网络接口层、网际层、运输层、应用层。\\
\textbf{五层原理体系结构}：教学用体系结构，自下而上为物理层、数据链路层、网络层、运输层、应用层。\\
\textbf{协议数据单元（PDU）}：对等层次之间传送的数据单位，如比特、帧、分组、报文段、报文。\\
\textbf{封装}：发送数据时在数据前面加上本层首部，逐层向下传递的过程。\\
\textbf{实体}：表示任何可以发送或接收信息的硬件或软件进程。\\
\textbf{客户／服务器（C/S）方式}：客户是服务请求方，服务器是服务提供方，客户主动发起请求。\\
\textbf{对等（P2P）方式}：通信双方地位对等，每个主机既可作为客户请求服务，也可作为服务器提供服务。\\
\textbf{RFC}：互联网正式标准的文档形式，由互联网工程任务组（IETF）等组织维护和发布。
\end{keyterms}

\begin{contrast}{OSI 七层 / TCP/IP 四层 / 五层（每行依次为：层数·优点·缺点或地位）}
\citem{OSI 七层}{7 层：物理、数据链路、网络、运输、会话、表示、应用}{层次划分清晰完整，理论严谨}{过于复杂，部分功能重复，实际未成为主流}
\citem{TCP/IP 四层}{4 层：网络接口层、网际层、运输层、应用层}{简单实用，已成为互联网事实标准}{网络接口层未细分，理论体系不够完整}
\citem{五层}{5 层：物理、数据链路、网络、运输、应用}{综合两者优点，便于教学和理解}{教学模型，工程上仍以 TCP/IP 为准}
\end{contrast}

\begin{contrast}{协议 / 服务（每行依次为：定义·作用范围·关系）}
\citem{协议}{同层对等实体之间进行通信的规则}{水平方向，同层之间}{协议的实现保证向上层提供服务}
\citem{服务}{下层向上层提供的功能}{垂直方向，相邻层之间}{服务由协议实现，上层使用下层服务}
\end{contrast}

\begin{contrast}{客户／服务器方式 / 对等（P2P）方式（每行依次为：通信角色·优点·典型应用）}
\citem{C/S}{客户主动请求，服务器被动响应并提供服务}{集中管理、易于维护和安全控制}{Web、电子邮件、FTP}
\citem{P2P}{通信双方地位对等，可同时作为客户和服务器}{可扩展性好、减轻服务器压力}{BT 下载、区块链、即时通信}
\end{contrast}

\qtypename{一、单项选择题}{10 题，每题 3 分，共 30 分}
\q 网络协议的三要素是（\quad）
\choices{A. 语法、语义和同步}{B. 硬件、软件和协议}
{C. 速率、带宽和时延}{D. 物理层、网络层和运输层}

\q 计算机网络体系结构的定义是（\quad）
\choices{A. 计算机网络各层及其协议的集合}{B. 网络中的硬件设备及其连接方式}
{C. 网络操作系统和应用程序的集合}{D. 网络拓扑结构与传输媒体的统称}

\q 在 OSI 参考模型中，负责为两台主机中进程之间提供端到端可靠传输的是（\quad）
\choicesline{A. 网络层}{B. 运输层}{C. 数据链路层}{D. 会话层}

\q OSI 参考模型自下而上第三层是（\quad）
\choicesline{A. 数据链路层}{B. 网络层}{C. 运输层}{D. 会话层}

\q TCP/IP 体系结构中，网际层的核心协议是（\quad）
\choicesline{A. IP}{B. TCP}{C. UDP}{D. HTTP}

\q 在五层原理体系结构中，数据链路层的协议数据单元是（\quad）
\choicesline{A. 比特}{B. 帧}{C. 分组}{D. 报文段}

\q 路由器工作于（\quad）
\choicesline{A. 物理层}{B. 数据链路层}{C. 网络层}{D. 应用层}

\q 应用层的协议数据单元通常称为（\quad）
\choicesline{A. 帧}{B. 分组}{C. 报文段}{D. 报文}

\q 在 OSI 参考模型中，负责数据表示、加密和压缩的是（\quad）
\choicesline{A. 会话层}{B. 表示层}{C. 应用层}{D. 运输层}

\q 在客户／服务器方式中，服务器是（\quad）
\choices{A. 服务请求方，被动等待}{B. 服务提供方，被动等待并响应请求}
{C. 与客户完全对等的节点}{D. 只负责数据转发的路由器}

\qtypename{二、多项选择题}{3 题，每题 4 分，共 12 分；每题有两个或两个以上正确项}
\q 下列属于 OSI 参考模型层次的有（\quad）
\choicesline{A. 会话层}{B. 表示层}{C. 运输层}{D. 网络接口层}

\q 下列属于 TCP/IP 体系结构层次的有（\quad）
\choicesline{A. 网际层}{B. 运输层}{C. 应用层}{D. 会话层}

\q 关于封装与解封装，下列说法正确的有（\quad）
\choices{A. 发送时每一层都在数据前面加上本层首部}{B. 接收端逐层剥去首部并向上交付}
{C. 物理层传输的是比特流}{D. 封装只发生在运输层}

\qtypename{三、判断题}{6 题，每题 2 分，共 12 分；正确的画“√”，错误的画“×”并写出改正后的表述}
\q 网络协议由语法、语义和同步三个要素组成。\quad（\quad）\quad 改正：\fixblank

\q OSI 参考模型分为 7 层，TCP/IP 体系结构分为 4 层，五层体系结构是在两者基础上形成的教学模型。\quad（\quad）\quad 改正：\fixblank

\q 服务是不同层对等实体之间通信的规则，协议是下层向上层提供的功能。\quad（\quad）\quad 改正：\fixblank

\q 在 OSI 参考模型中，网络层负责实现端到端的可靠传输。\quad（\quad）\quad 改正：\fixblank

\q 数据链路层的协议数据单元是报文段。\quad（\quad）\quad 改正：\fixblank

\q 客户／服务器方式中，客户主动发起请求，服务器被动等待并响应请求。\quad（\quad）\quad 改正：\fixblank

\qtypename{四、填空题}{6 空，每空 1 分，共 6 分}
\q 网络协议的三要素是语法、语义和\fillblanks[2.0cm]。

\q OSI 开放系统互连参考模型共分为\fillblanks[1.2cm]层。

\q TCP/IP 体系结构中最上面一层是\fillblanks[2.0cm]。

\q 五层原理体系结构中，网络层的协议数据单元是\fillblanks[2.0cm]。

\q 数据链路层的协议数据单元是\fillblanks[1.6cm]。

\q 互联网的正式标准都以\fillblanks[2.0cm]文档的形式发表。

\qtypename{五、简答题}{2 题，每题 10 分，共 20 分}
\q 简述网络协议的三要素，并说明计算机网络采用分层体系结构的好处。\anslines{2.4cm}

\q 写出 OSI 七层和 TCP/IP 四层体系结构的各层名称，说明二者的区别与联系，并简述五层原理体系结构各层的主要功能。\anslines{2.4cm}

\newpage
\qtypename{六、综合题}{1 题，共 20 分}
\q 用户用浏览器访问某农业信息网站。请结合数据从发送端到接收端的传输过程回答：（1）写出应用层、运输层、网络层、数据链路层和物理层分别使用的典型协议（HTTP、TCP、IP、以太网）及对应的协议数据单元名称；（2）说明发送端逐层封装和接收端逐层解封装的过程；（3）说明交换机、路由器分别工作在哪个层次、转发数据时依据什么信息；（4）用表格或文字说明 OSI 七层与 TCP/IP 四层、五层体系结构的对应关系，并指出互联网实际使用的是哪一种体系结构。\anslines{5cm}

\end{document}
